% Pursuing Stacks (cotes 134-2 to 134-8): archive cote 134-8, site page 3 (PDF page 4), typescript page 557.
% typescript. Draft compared: ps7.tex:86-153.
% AI-transcribed by Claude Opus 5.5 (claude-opus-5-5, adaptive thinking, effort high), Message Batches API, batch ps4 (msgbatch_012Aoc1NuskATPsyBeKSYK7J), 2026-11-12; not yet checked by eye.
\scanpage{134-8/3}
\mspage{557}

It seems to me there has'nt been a single thing in mathematics,  \uline{inclu-}%
\uline{ding motives},  which has exerted such a fascination upon me.  It will be \margin{\unsure{it's very close anyhow}}
hard I am affraid to carry on a seminar on such a stuff,  and go on and
carry to their (hopefully happy!) end those ponderings and notes on
homotopical algebra so-called  -  and a somewhat crazy one too at times,
I am affraid!  We'll see what comes out of all this !  T\unsure{\over{n}{h}}e night after
the seminar session at any rate,  and already during the two hours dri-%
\xout{\unsure{At any rate I decide}}
ve home,  \xout{the head} \deleted{the stuff} \inserted{the Teichmüller stuff} was brewing anew in my head.  Maybe it
would have gone on for days,  but \deleted{next} \inserted{the very} day already several things came \margin{\unsure{$\sim$}}
up demanding special attention,
% NOTE: a white box hides the rest of this line and the beginning of the next line; the next line is hidden as far as the box reaches

26.8.

134.  It is time to come back to the ``review'' of linearization of
(ordinary) homotopy type\hand{s},  and the homology and cohomology formalism
in the context of the modelizer (Cat),  which has been pushed aside now
for over two months,  for the sake of that endless digression on schema-%
tization of homotopy types.\hand{\unsure{*}}  The strong tie between these two strains
of reflection has been the equal importance in both of the linearization
process.  Technically speaking though,  linearization\hand{,} as finally handled
in the previous chapter (via the so-called ``integrators'')\hand{,} looks pretty
much different from the similar operation in the schematic set-up,  due
\hand{(}at least partly\hand{)} to the choice we made of models for expressing schematic
homotopy types,  namely taking semisimplicial scheme-like objects\over{,}{;}
which means,  notably,  working with \hand{$\Delta$} rather than more general test
categories,  and relying heavily on the Kan-Dold-Puppe relationship.
In the discrete set-up,  it turned out (somewhat inexpectedly) that the
latter can be replaced,  when working with more general categories \inserted{A} than
\hand{$\Delta$} (which need not even be test categories or anything of the kind),  by
the L\hand{$p$}$^{ab}_{!}$ operation (when \hand{$p$} : A $\longrightarrow$ e is the map from A to the final ob- \margin{\textcircled{p}}
ject e in (Cat)),  computable  in terms of a choice of an ``integrator''
for A.  More specifically,  recalling that
\[ (1) \qquad p^{ab}_{!} : A^{\wedge}_{ab} \longrightarrow (Ab) \]
is defined as the left adjoint of the  inverse image functor
\[ p^{*}: \ (Ab) \longrightarrow A^{\wedge}_{ab} \quad , \]
associating to every abelian group the corresponding ``constant'' abelian
\margin{3}
